\caption{\textbf{Cortical--subcortical separation in functional-cluster composition space.}
For each of 198 Beryl regions ($\geq$20 neurons; 68 cortical, 130 subcortical; all trials, 47,378 neurons), the fraction of its neurons in each of the 25 k-means clusters of Figs.~2--3 (composition vector).
\textbf{a} Cell-count-weighted Fisher criterion $D$ between cortical and subcortical composition vectors (vertical line, $D=0.62$) and its null distribution from 1,000 region-label shuffles ($0.20\pm0.05$; $p=0.001$, the smallest possible).
\textbf{b} PERMANOVA on Euclidean distances between composition vectors: significant but small group effect (pseudo-$F=10.2$, $R^2=0.050$, $p=0.001$).
\textbf{c} Specialization ($1-H(\mathbf{p})/\log K$) does not differ between cortical and subcortical regions (Mann--Whitney $p=0.80$).
\textbf{d} PCA of the composition vectors (fit with rows weighted by $\sqrt{n_\mathrm{cells}}$, axes centred on the weighted mean; PC1 24\%, PC2 18\% of the variance). Regions in Beryl colours; circles cortical, squares subcortical; dashed ellipses, 1\,s.d.; the ten largest regions labelled. Cortical and subcortical regions are shifted along PC1 but overlap broadly.}
\label{fig:cortical_subcortical_specialization}
